\subsection{Arithmetic} \label{sec:math-task}

Modern LMs have been shown to possess basic numerical reasoning abilities~\citep{lewkowycz2022solving}, with \citet{gpt3} even reporting near-perfect GPT-3 accuracy for two-digit additions. On the other hand, \citet{razeghi-etal-2022-impact} find that LMs perform significantly better on operations involving numbers that occur more frequently in the pretraining data, and \citet{li-etal-2023-bert} show that symbol replacement affects the mathematical ability of BERT~\citep{devlin2019bert}-like models; both findings point to overfitting and memorization effects. We consider the same two-digit addition task, the simplest arithmetic task in \citet{gpt3}, but inspect a model's accuracy in different bases. We use base-8, 9, 11, and 16 as the counterfactual setup which are natural generalizations to base-10 arithmetic. These bases were chosen to control for task difficulty (see \S\ref{sec:underestimation} for a discussion) and also to test for how relatively uncommon (9 \& 11) and common (8 \& 16) bases affect performance (see \S\ref{sec:analysis-world-commonness} for an analysis).
To ensure the model understands the different bases, the CCC evaluates the successor relation under each base.
